\subsubsection{Các mô hình hàm cơ sở phi tuyến}

Tất cả các mô hình đều tuân theo dạng tổng quát:

\begin{equation}
f(\mathbf{x}, \mathbf{w}) = \mathbf{w}^T \Phi(\mathbf{x})
\end{equation}

trong đó $\Phi(\mathbf{x})$ được xây dựng từ các hệ cơ sở khác nhau.

\paragraph{Mô hình đa thức (Polynomial Basis):}
Mô hình đa thức mở rộng vector đầu vào thành các lũy thừa phi tuyến:

\begin{equation}
\Phi(\mathbf{x}) = [1, x, x^2, \dots, x^p]
\end{equation}

Đây là dạng đơn giản nhất, giúp mô hình xấp xỉ các quan hệ phi tuyến mượt.

\textbf{Ưu điểm:}
\begin{itemize}
    \item Dễ cài đặt
    \item Diễn giải tốt
\end{itemize}

\textbf{Nhược điểm:}
\begin{itemize}
    \item Dễ overfitting khi bậc lớn
    \item Không ổn định với dữ liệu lớn
\end{itemize}

\paragraph{Mô hình RBF (Radial Basis Function):}
Mô hình RBF sử dụng các hàm Gaussian làm cơ sở:

\begin{equation}
\phi_j(\mathbf{x}) =
\exp\left(-\frac{\|\mathbf{x} - \mathbf{c}_j\|^2}{2\sigma^2}\right)
\end{equation}

Trong đó $\mathbf{c}_j$ là các tâm và $\sigma$ điều khiển độ lan tỏa.

Mô hình tập trung vào khoảng cách trong không gian đầu vào, giúp biểu diễn tốt cấu trúc phi tuyến cục bộ.

\textbf{Ưu điểm:}
\begin{itemize}
    \item Xấp xỉ tốt dữ liệu phi tuyến
    \item Linh hoạt theo phân bố dữ liệu
\end{itemize}

\paragraph{Mô hình Fourier Basis:}
Mô hình Fourier biểu diễn dữ liệu bằng các hàm sin/cos:

\begin{equation}
\Phi(\mathbf{x}) =
[\sin(\omega_1 x), \cos(\omega_1 x), \dots, \sin(\omega_k x), \cos(\omega_k x)]
\end{equation}

Biến đổi dữ liệu sang miền tần số, phù hợp với dữ liệu có tính chu kỳ.

\textbf{Ưu điểm:}
\begin{itemize}
    \item Trực giao tốt
    \item Ổn định số học cao
\end{itemize}

\paragraph{So sánh tổng quát:}
Ba nhóm mô hình có đặc điểm khác nhau:

\begin{itemize}
    \item Polynomial: đơn giản nhưng dễ overfit
    \item RBF: mạnh với dữ liệu cục bộ
    \item Fourier: tốt cho dữ liệu tuần hoàn
\end{itemize}
